% Imporditud: tools/import_eksperimendid.py
% Allikas: Eesti füüsikaolümpiaadi eksperimendi ülesannete
%   kogu (Rael Kalda, 2023), ülesanne Ü105, lahendus L105.
\setAuthor{EFO žürii}
\setRound{piirkonnavoor}
\setYear{2012}
\setNumber{G E2}
\setDifficulty{3}
\setTopic{dünaamika}
\setVahendid{kuuekandiline pliiats, niit, koormis, stopper, statiiv kinnitusklambriga}
\setPunktid{12}
\setAllikas{Eksperimendi kogu Ü105, lahendus lk 80}
\setLahendusseis{toores}

\prob{Pliiats}
Määrake pliiatsi ristlõike ümbermõõt. Eeldage, et raskuskiirendus g = 9,8 m/s$^2$. Mõõtemääramatust ei ole vaja leida.

\hint

\solu
Mähime ümber pliiatsi N niidikeerdu, nt N = 10 (mõistlikult suur arv keerde: 1 p) ning kasutame sellele tiirude arvule vastavat niidi osa matemaatilise pendli ehitamiseks (2 p) (vähendades niidi pikkust selle vahemaa võrra, mis eraldab niidi ja koormise ühenduspunkti koormise masskeskmest (1 p)). Mõõdame stopperiga võnkeperioodi T (3 p) kasutades mõõtetäpsuse huvides suurt võngete arvu (1 p), nt M = 20. Seosest T = 2$\pi$ l/g (2 p) saame pliiatsi ümbermõõdu p = l/N = T 2g/4$\pi$2N (1 p). Numbriline vastus (1 p).
\probend
